
\documentclass[10pt]{article}
\usepackage[utf8]{inputenc}
\usepackage[T1]{fontenc}
\usepackage{amsmath, amssymb, xcolor, geometry, enumitem, multicol, tcolorbox}

\definecolor{mainblue}{RGB}{0, 102, 204}
\geometry{a4paper, margin=1.2cm, top=1cm, bottom=3cm, footskip=1.2cm}

\begin{document}
\enlargethispage{2.5cm}

% Modern Header
\begin{tcolorbox}[colback=mainblue!5, colframe=mainblue, sharp corners, boxrule=0.5pt]
    \small \textbf{Unit:} undefined \hfill \textbf{Name:} \rule{3cm}{0.4pt} \\
    \small \textbf{Topic:} Variables and Expressions / Order of Operations \hfill \textbf{Date:} \rule{2cm}{0.4pt} \\
    \centering \textbf{\large Challenge (Level 0)}
\end{tcolorbox}

\vspace{0.2cm}

% MCQ Section
\begin{multicols}{2}
\begin{enumerate}[label=\textbf{Q\arabic*.}, leftmargin=*, itemsep=6pt, parsep=0pt]

  \item \small What is the variable in the expression $2x + 3$?
  \begin{enumerate}[label=(\Alph*), leftmargin=0.6cm, nosep, itemsep=2pt]
    \item 2
    \item 3
    \item x
    \item 2x
    \item None of the above
  \end{enumerate}

  \item \small Identify the constant term in the expression $4y - 2$.
  \begin{enumerate}[label=(\Alph*), leftmargin=0.6cm, nosep, itemsep=2pt]
    \item 4y
    \item 2
    \item 4
    \item -2
    \item -4
  \end{enumerate}

  \item \small What is the coefficient of $x$ in the expression $5x - 1$?
  \begin{enumerate}[label=(\Alph*), leftmargin=0.6cm, nosep, itemsep=2pt]
    \item 1
    \item -1
    \item 5
    \item -5
    \item None of the above
  \end{enumerate}

  \item \small What operation is denoted by the symbol $+$ in the expression $2 + 3$?
  \begin{enumerate}[label=(\Alph*), leftmargin=0.6cm, nosep, itemsep=2pt]
    \item Subtraction
    \item Addition
    \item Multiplication
    \item Division
    \item None of the above
  \end{enumerate}

  \item \small Identify the operator in the expression $7 - 2$.
  \begin{enumerate}[label=(\Alph*), leftmargin=0.6cm, nosep, itemsep=2pt]
    \item 7
    \item 2
    \item -
    \item +
    \item *
  \end{enumerate}

  \item \small What is the term $3x$ an example of in the expression $3x + 2$?
  \begin{enumerate}[label=(\Alph*), leftmargin=0.6cm, nosep, itemsep=2pt]
    \item Constant term
    \item Variable term
    \item Coefficient
    \item Operator
    \item None of the above
  \end{enumerate}

  \item \small In the expression $9 - 4$, what is being subtracted?
  \begin{enumerate}[label=(\Alph*), leftmargin=0.6cm, nosep, itemsep=2pt]
    \item 9
    \item 4
    \item -
    \item +
    \item *
  \end{enumerate}

  \item \small What is the purpose of the order of operations in an expression like $(2 + 3) \cdot 4$?
  \begin{enumerate}[label=(\Alph*), leftmargin=0.6cm, nosep, itemsep=2pt]
    \item To simplify the expression
    \item To evaluate the expression
    \item To identify the variables
    \item To determine the coefficients
    \item None of the above
  \end{enumerate}

\end{enumerate}
\end{multicols}

\vspace{-0.2cm}
\hrule
\vspace{0.2cm}
\textbf{\small Open Ended Questions (FRQ)}
\begin{enumerate}[label=\textbf{Q\arabic*.}, start=9, itemsep=5pt, topsep=0pt]

  \item \small Draw a diagram to represent the expression $x + 2$. Label the variable and the constant. \vspace{2cm}

  \item \small Explain the concept of the order of operations in your own words. Provide an example of how it is used in a simple expression like $3 + 2 \cdot 4$. \vspace{4cm}

\end{enumerate}

\vfill

% Chef's Tips Footer
\begin{tcolorbox}[colback=blue!5, colframe=mainblue!50, title=\small \textbf{Chef's Tips}, fonttitle=\bfseries, bottom=1mm]
\begin{itemize}[noitemsep, topsep=2pt, leftmargin=0.5cm]
  \item \scriptsize Focus on basic conceptual definitions and single-step identification for Level 0.\item \scriptsize Ensure strict adherence to the order of operations when evaluating expressions.\item \scriptsize Use visual aids and diagrams to help students understand the concepts of variables, constants, and operators.
\end{itemize}
\end{tcolorbox}

\end{document}
  